\documentclass[10pt,a4paper]{article}
\pdfinfoomitdate=1
\pdftrailerid{}
\pdfsuppressptexinfo=-1
\usepackage[T1]{fontenc}
\usepackage[utf8]{inputenc}
\usepackage{mathpazo}
\IfFileExists{beramono.sty}{\usepackage[scaled=0.86]{beramono}}{}
\usepackage[a4paper,left=2.22cm,right=2.22cm,top=2.4cm,bottom=2.6cm]{geometry}
\usepackage{amsmath,amssymb,amsthm}
\usepackage{array,longtable}
\usepackage[table]{xcolor}
\usepackage{microtype}
\DeclareUnicodeCharacter{211D}{\ensuremath{\mathbb{R}}}
\DeclareUnicodeCharacter{2264}{\ensuremath{\le}}
\DeclareUnicodeCharacter{2200}{\ensuremath{\forall}}
\DeclareUnicodeCharacter{2228}{\ensuremath{\lor}}
\usepackage[hidelinks,
  pdftitle={Positive roots of sums of cubes of binomials: a readable account of eighteen Lean-checked statements, with one new case},
  pdfauthor={},pdfsubject={},pdfkeywords={}]{hyperref}

% Layout in the manner of the project's first paper: left-aligned bold title, a lead paragraph, numbered bold
% headings with a hairline, ragged-right text without indentation, ruled tables with a shaded header row.
\makeatletter
\renewcommand\normalsize{\@setfontsize\normalsize{10.6}{15.4}%
  \abovedisplayskip 8pt plus 2pt minus 4pt \abovedisplayshortskip 2pt plus 3pt
  \belowdisplayshortskip 5pt plus 3pt minus 3pt \belowdisplayskip\abovedisplayskip \let\@listi\@listI}
\normalsize
\renewcommand\small{\@setfontsize\small{9.2}{13}}
\renewcommand\footnotesize{\@setfontsize\footnotesize{8.6}{12.2}}
\renewcommand\section{\@startsection{section}{1}{\z@}{1.7em plus .3em minus .2em}{0.25em}%
  {\normalfont\bfseries\fontsize{13.5}{17}\selectfont\raggedright}}
\renewcommand\subsection{\@startsection{subsection}{2}{\z@}{1.3em plus .2em minus .2em}{0.35em}%
  {\normalfont\bfseries\fontsize{11.6}{15}\selectfont\raggedright}}
\renewcommand\@seccntformat[1]{\csname the#1\endcsname.\ }
\let\note@section\section
\RenewDocumentCommand{\section}{s o m}{%
  \IfBooleanTF{#1}{\note@section*{#3}}{\IfNoValueTF{#2}{\note@section{#3}}{\note@section[#2]{#3}}}%
  \par\nobreak\vspace{-0.15em}{\color{black!22}\hrule height 0.5pt}\nobreak\vspace{0.55em}\@afterheading}
\makeatother
\definecolor{headgray}{gray}{0.93}
\arrayrulecolor{black!25}
\setlength{\arrayrulewidth}{0.5pt}
\renewcommand{\arraystretch}{1.3}
\setlength{\parindent}{0pt}
\setlength{\parskip}{0.62em}
\pagestyle{empty}
\AtBeginDocument{\raggedright}
\emergencystretch=1em

\theoremstyle{definition}
\newtheorem{theorem}{Theorem}
\newtheorem{lemma}[theorem]{Lemma}
\newtheorem{proposition}[theorem]{Proposition}
\newtheorem{corollary}[theorem]{Corollary}
\newtheorem{example}[theorem]{Example}

\newcommand{\R}{\mathbb{R}}
\newcommand{\lean}[1]{{\ttfamily\small\renewcommand{\_}{\textunderscore\allowbreak}#1}}
\newcommand{\row}[1]{L-gct-#1}
\newcommand{\eps}{\varepsilon}
\newcolumntype{L}[1]{>{\raggedright\arraybackslash}p{#1}}

\begin{document}
\thispagestyle{empty}

{\raggedright\bfseries\fontsize{17}{22}\selectfont Positive roots of sums of cubes of binomials: a readable
account of eighteen Lean-checked statements, with one new case\par}
\vspace{0.9em}

This document is a human-readable account of eighteen statements about the roots of small sums of cubes and of
products of sparse real polynomials, each checked in Lean~4 with Mathlib by contributors to the PEqualNP project,
together with one new example and its consequence. The Lean files are public at
\url{https://github.com/open-goat/pequalnp-proofs}, under \lean{lean/PEqualNP/Outside/}: one file per statement,
named in Table~\ref{tab:rows}. The text says for each statement whether it is checked in Lean, proved here on
paper, or the result of a computation. The prose itself is not machine-checked and has not been read by a human
mathematician.

A nonzero sum of $k$ cubes of real binomials has at most $4k$ monomials, hence at most $4k-1$ positive roots by
Descartes' rule of signs. Contributors to the project have given examples, checked in Lean, with $4k-1$ positive
roots for $k=3$, $5$ and $6$, and with $14$ for $k=4$. This note gives an example with $15$ positive roots for
$k=4$, also checked in Lean, and an elementary gluing lemma. Together with the earlier examples they show that the
largest number of positive roots of a nonzero sum of $k$ cubes of real binomials is $4k-1$ for every $k\ge3$; it is
$1$ for $k=1$ and $3$ for $k=2$. For the order of a positive root of a sum of three cubes, the project has proved
the bound $9$ for polynomials and an example of order $7$ is known. The note observes that the same proof gives the
bound $9$ for real exponents, records an identity with real exponents in which the order is exactly $9$, and
reports an exact computation showing that no example with integer exponents up to $32$ has order $8$ or more. It
also describes the related checked statements on $fg+1$ and $pq-rs$. It says nothing about the truth of the real
$\tau$-conjecture.

26 November 2026.

\section{Introduction}\label{sec:intro}

A \emph{$t$-nomial} is a real polynomial with at most $t$ monomials; binomials and trinomials are the cases $t=2,3$.
By Descartes' rule of signs a nonzero $t$-nomial has at most $t-1$ positive roots, and for every choice of $t$
exponents some choice of coefficients attains this. The question behind this note is what remains of the bound when
the coefficients are not free. The simplest structured families are $\sum_{i=1}^k g_i^3$ with binomials $g_i$,
where $4k$ coefficients depend on $2k$ parameters, and $fg+1$ with $t$-nomials $f,g$. Questions of this kind come
from Koiran's real $\tau$-conjecture, which predicts that sums of products of sparse polynomials have few real
roots (Section~\ref{sec:context}).

The PEqualNP project (\url{https://pequalnp.com}) maintains a public repository of Lean~4 proofs,
\url{https://github.com/open-goat/pequalnp-proofs}. In November 2026 it posted forty statements on small cases of
this question, as twenty pairs consisting of a bound and its negation, and outside contributors proved one statement
of eighteen of the pairs. Table~\ref{tab:rows} lists these eighteen statements. Each is the last theorem of a file
\lean{lean/PEqualNP/Outside/L\_gct\_N.lean}. The column ``statement'' gives what was posted and proved. The column
``file'' gives the number of roots that the file itself exhibits, which is often larger. The column $D$ is the upper
bound given by the rule of signs. The last column names the account that submitted the accepted pull request.

\begin{table}[ht]
\centering\footnotesize
\setlength{\tabcolsep}{4pt}
\begin{tabular}{|l|L{5.0cm}|c|c|c|L{3.9cm}|}
\hline
\rowcolor{headgray}\textbf{name} & \textbf{family} & \textbf{statement} & \textbf{file} & $\boldsymbol{D}$ & \textbf{submitted by (pull request)}\\\hline
\row{241} & $3$ cubes of binomials, positive roots & $>9$ & $11$ & $11$ & XXXLUMI (\#14351)\\\hline
\row{251} & $4$ cubes of binomials, positive roots & $>13$ & $14$ & $15$ & k-sandhu (\#14687)\\\hline
\row{253} & $5$ cubes of binomials, positive roots & $>17$ & $19$ & $19$ & junyeopYim (\#15826)\\\hline
\row{255} & $6$ cubes of binomials, positive roots & $>21$ & $23$ & $23$ & k-sandhu (\#14652)\\\hline
\row{257} & $3$ cubes of binomials, real roots & $>19$ & $22$ & $23$ & k-sandhu (\#14650)\\\hline
\row{243} & $3$ cubes of trinomials, positive roots & $>14$ & $15$ & $29$ & k-sandhu (\#14699)\\\hline
\row{245} & $3$ cubes of $4$-nomials, positive roots & $>21$ & $23$ & $59$ & k-sandhu (\#14662)\\\hline
\row{247} & $3$ cubes of $5$-nomials, positive roots & $>26$ & $27$ & $104$ & k-sandhu (\#14717)\\\hline
\row{249} & $3$ cubes of $6$-nomials, positive roots & $>33$ & $35$ & $167$ & k-sandhu (\#14669)\\\hline
\row{259} & $3$ cubes of trinomials, real roots & $>29$ & $30$ & $59$ & k-sandhu (\#14706)\\\hline
\row{234} & $3$ cubes of binomials, order of a positive root &
\multicolumn{2}{c|}{$\le 9$} & $11$ & Aargau (\#23576)\\\hline
\row{261} & $fg+1$, trinomials, positive roots & $>7$ & $8$ & $9$ & hedanbaomi (\#13871)\\\hline
\row{263} & $fg+1$, $4$-nomials, positive roots & $>10$ & $11$ & $16$ & k-sandhu (\#14768)\\\hline
\row{265} & $fg+1$, $5$-nomials, positive roots & $>11$ & $12$ & $25$ & k-sandhu (\#14655)\\\hline
\row{267} & $fg+1$, trinomials, real roots & $>7$ & $8$ & $19$ & BastianMIllan (\#13649)\\\hline
\row{269} & $fg+1$, $4$-nomials, real roots & $>12$ & $13$ & $33$ & blarneys (\#13524)\\\hline
\row{271} & $fg+1$, $5$-nomials, real roots & $>14$ & $15$ & $51$ & gcq (\#13516)\\\hline
\row{273} & $pq-rs$, $4$-nomials, positive roots & $>22$ & $23$ & $31$ & SemyonOsadchii (\#16079)\\\hline
\end{tabular}
\caption{The eighteen checked statements. Coefficients are real. The table counts distinct roots, except for
\row{234}, which bounds the order of a single root.}
\label{tab:rows}
\end{table}

Throughout, bounds on the number or the order of roots concern functions that are not identically zero. The file of
\row{234} proves an upper bound for the order of a root of a polynomial. The other files certify lower bounds for
the number of roots of explicit examples. When an example attains the bound of the rule of signs, that bound also
gives the exact number of roots and their simplicity. The general statement on the number of roots,
Theorem~\ref{thm:main}, uses in addition the gluing lemma.

The multiplicity question was also studied in Discussion \#21999 of the repository. The following results from it
are used below. They are due to maxkieff: the polynomial
$2250(X^3-X^2)^3-125(X^4-X)^3+9(X^5-1)^3$ has a root of order exactly $7$ at $1$; if all three binomials vanish at
the root, the order is at most $7$; and, in that case, the conditions that order $8$ imposes on the exponents, which
are recalled before Proposition~\ref{prop:family}. The discussion also contains a search by maxkieff excluding
orders $10$ and $11$ for small exponents, and an argument of silentoplayz excluding order $11$ when the sum has
twelve distinct exponents, and order $10$ when it has eleven, provided that three consecutive exponents of one cube
are not shared with another cube and no other exponent of the sum lies between them. Both were overtaken by
\row{234}, which was proved a few hours later.

Beyond the eighteen files and that discussion, this note contains the following.
\begin{enumerate}\setlength{\itemsep}{0pt}
\item[(a)] A sum of four cubes of binomials with $15$ positive roots (Theorem~\ref{thm:four}). The lower bound $15$ is
checked in Lean.
\item[(b)] A gluing lemma (Lemma~\ref{lem:glue}) and the resulting exact values for all $k$
(Theorem~\ref{thm:main}, Corollary~\ref{cor:real}). These are proved here and are not formalized.
\item[(c)] The remark that the proof of \row{234} gives the bound $9$ also for real exponents
(Proposition~\ref{prop:real}), and an identity showing that $9$ is attained there (Proposition~\ref{prop:nine}),
together with the next two moments of the family studied by maxkieff (Proposition~\ref{prop:family}).
\item[(d)] An exact computation: no nonzero sum of three cubes of binomials with exponents in $[0,32]$ has a
positive root of order $8$ or more (Proposition~\ref{prop:search}).
\end{enumerate}
All eighteen files of Table~\ref{tab:rows} were recompiled against Mathlib alone, and the statements quoted here were
compared with the Lean declarations (Appendix~\ref{sec:lean}).

\section{Preliminaries}\label{sec:prelim}

Let $\mathcal{S}_k$ be the set of real polynomials $\sum_{i=1}^k g_i^3$ with binomials $g_i$. The Lean statements
are written for $\sum_i s_i g_i^3$ with $s_i\in\{1,-1\}$, which is the same set because $-g^3=(-g)^3$; real
coefficients $c_ig_i^3$ also give the same set. Let $N(k)$ be the largest number of distinct positive roots of a
nonzero polynomial in $\mathcal{S}_k$.

Section~\ref{sec:mult} also uses \emph{signomials} $\sum_j c_jx^{e_j}$ with real exponents, regarded as functions on
$(0,\infty)$. The two lemmas below hold for them with the same proofs.

\begin{lemma}[Descartes' rule of signs]\label{lem:descartes}
A nonzero real signomial with $m$ monomials has at most $m-1$ positive roots counted with multiplicity.
\end{lemma}

\begin{lemma}\label{lem:hajos}
A nonzero real signomial with $m$ monomials has no positive root of order $m$ or more.
\end{lemma}

Lemma~\ref{lem:descartes} is classical; see \cite{AlbouyFu}, which also recalls that the bound is attained for free
coefficients. Lemma~\ref{lem:hajos} is a special case. It is usually attributed to Haj\'os \cite{Hajos}; printed
proofs are in \cite[Proposition~3.2]{Lenstra} and \cite[Lemma~3.9]{KoiranSkomra}, and a proof by a Vandermonde
determinant is formalized in the file of \row{234}. Since a polynomial in $\mathcal{S}_k$ has at most $4k$
monomials,
\begin{equation}\label{eq:upper}
N(k)\le4k-1 .
\end{equation}

A function $f$ on $(0,\infty)$ \emph{alternates} at the points $x_0<x_1<\dots<x_n$ if $f(x_i)\ne0$ for all $i$ and
$f(x_{i-1})f(x_i)<0$ for $i=1,\dots,n$. A continuous function that alternates at $n+1$ points has at least $n$
distinct zeros. Alternation is used to prove the root counts of Sections \ref{sec:three} and \ref{sec:more}.

\begin{proposition}\label{prop:two}
$N(1)=1$ and $N(2)=3$.
\end{proposition}
\begin{proof}
A nonzero binomial has at most one positive root, and $(x-1)^3$ has one. Since $y\mapsto y^3$ is injective on $\R$,
$g_1(x)^3+g_2(x)^3=0$ exactly when $g_1(x)+g_2(x)=0$. If the sum of cubes is not the zero polynomial, $g_1+g_2$ is
a nonzero $4$-nomial and has at most three positive roots. With $g_1=x^3-6x^2$ and $g_2=11x-6$ the roots $1,2,3$
are attained.
\end{proof}

Dutta \cite[Theorem~10]{Dutta} uses the same observation for sparse polynomials in general: the number of real roots
of a nonzero $c_1f_1^3+c_2f_2^3$ is at most a constant times the number of monomials occurring in $f_1$ or $f_2$.
For two cubes of binomials the expansion has up to eight monomials, so the rule of signs alone would allow seven
positive roots.

\section{Three cubes}\label{sec:three}

\begin{theorem}[\row{241}, XXXLUMI]\label{thm:three}
Let $\eps=39/250000$, $c=1678911/20000000$ and
\[
g_0=x^{6402}-\eps\,x^{12802},\qquad g_1=\eps-x^{6400},\qquad g_2=c\,\bigl(x^{8449}-x^{4353}\bigr).
\]
The polynomial $P=g_0^3+g_1^3+g_2^3$ has exactly $11$ positive roots, all simple. Hence $N(3)=11$.
\end{theorem}
\begin{proof}
One has $x^{12802}g_0(1/x)=-g_1(x)$, $x^{12802}g_1(1/x)=-g_0(x)$ and $x^{12802}g_2(1/x)=-g_2(x)$, hence
\begin{equation}\label{eq:anti}
x^{38406}\,P(1/x)=-P(x)\qquad(x>0).
\end{equation}
Next, write $q=x^{256}$. Since $6400=256\cdot25$, $6402=2+256\cdot25$, $12802=2+256\cdot50$,
$4353=1+256\cdot17$ and $8449=1+256\cdot33$,
\[
g_0=x^2q^{25}(1-\eps q^{25}),\qquad g_1=\eps-q^{25},\qquad g_2=c\,x\,q^{17}(q^{16}-1),
\]
so that $P(x)=A(q)x^6+B(q)x^3+C(q)$ with $A=q^{75}(1-\eps q^{25})^3$, $B=c^3q^{51}(q^{16}-1)^3$ and
$C=(\eps-q^{25})^3$. Table~\ref{tab:three} gives the three terms at the six points $x_j=q_j^{1/256}$.

\begin{table}[ht]\centering\small
\begin{tabular}{|l|r|r|r|r|}
\hline
\rowcolor{headgray}$q_j$ & $A(q_j)\,x_j^6$ & $B(q_j)\,x_j^3$ & $C(q_j)$ & $P(x_j)$\\\hline
$26/25$ & $18.9390296$ & $0.0029101$ & $-18.9419289$ & $1.08\cdot10^{-5}$\\\hline
$27/25$ & $320.754188$ & $0.428191$ & $-321.182581$ & $-2.01\cdot10^{-4}$\\\hline
$17/15$ & $11842.5583$ & $92.2812$ & $-11934.8174$ & $2.21\cdot10^{-2}$\\\hline
$13/11$ & $268964.844$ & $7280.02$ & $-276246.583$ & $-1.72$\\\hline
$7/5$ & $2.43\cdot10^{9}$ & $1.72\cdot10^{11}$ & $-9.11\cdot10^{10}$ & $8.29\cdot10^{10}$\\\hline
$2$ & $-5.50\cdot10^{33}$ & $3.78\cdot10^{26}$ & $-3.78\cdot10^{22}$ & $-5.50\cdot10^{33}$\\\hline
\end{tabular}
\caption{The values of $P$ at $x_j=q_j^{1/256}$, rounded.}\label{tab:three}
\end{table}

The signs of $P(x_1),\dots,P(x_6)$ are $+,-,+,-,+,-$. The Lean file proves this with rational arithmetic: for
$j\le5$ the coefficients $A(q_j)$ and $B(q_j)$ are nonnegative, so $A(q_j)x^6+B(q_j)x^3+C(q_j)$ is nondecreasing in
$x\ge 0$, and it suffices to evaluate it at a rational lower or upper bound for $x_j$ obtained by eight successive
squarings; for $j=6$ the bounds $1<x_6<2$ suffice. By \eqref{eq:anti}, $P(1/x_j)$ and $P(x_j)$ have opposite
signs. Therefore $P$ alternates at the twelve points
\[
1/x_6<\dots<1/x_1<x_1<\dots<x_6 ,
\]
and has at least eleven positive roots. It has twelve monomials, so by Lemma~\ref{lem:descartes} it has at most
eleven positive roots counted with multiplicity.
\end{proof}

In increasing order the exponents of $P$ are $0$, $6400$, $12800$, $13059$, $17155$, $19200$, $19206$, $21251$,
$25347$, $25606$, $32006$, $38406$; they come from $g_1,g_1,g_1,g_2,g_2,g_1,g_0,g_2,g_2,g_0,g_0,g_0$, and the
coefficients alternate in sign along this list.

\section{Any number of cubes}\label{sec:more}

\begin{lemma}[gluing]\label{lem:glue}
Let $P\in\mathcal{S}_{k_1}$ alternate at $p+1$ points of $(0,\infty)$ and let $Q\in\mathcal{S}_{k_2}$ alternate at
$q+1$ points of $(0,\infty)$, where $p,q\ge0$. Then some $R\in\mathcal{S}_{k_1+k_2}$ alternates at $p+q+2$ points
of $(0,\infty)$.
\end{lemma}
\begin{proof}
Let $x_0<\dots<x_p$ and $y_0<\dots<y_q$ be the points for $P$ and $Q$, let $d$ be the degree of $P$, and fix an
integer $n$ with $3n\ge d+2$. Choose $\sigma\in\{1,-1\}$ such that $\sigma Q(y_0)$ and $P(x_p)$ have opposite
signs. For a real number $\lambda\ge1$ put
\[
R_\lambda(x)=P(x)+\sigma\,\lambda^{\,d+1-3n}\,x^{3n}\,Q(x/\lambda).
\]
If $Q=\sum_j h_j^3$ with binomials $h_j$, the second term is the sum of the cubes of the binomials
$\sigma\lambda^{(d+1-3n)/3}x^{n}h_j(x/\lambda)$, so $R_\lambda\in\mathcal{S}_{k_1+k_2}$.

At a point $x_i$ the second term is at most $\lambda^{-1}x_p^{3n}\max_{[0,x_p]}|Q|$ in absolute value, because
$d+1-3n\le-1$ and $0<x_i/\lambda\le x_p$. This tends to $0$ as $\lambda\to\infty$, so $R_\lambda(x_i)$ has the sign
of $P(x_i)$ for all large $\lambda$.

At a point $\lambda y_j$ the second term equals $\sigma\lambda^{d+1}y_j^{3n}Q(y_j)$, while
$|P(\lambda y_j)|\le\lambda^{d}\,c_P$ with a constant $c_P$ depending only on $P$ and $y_q$. So
$R_\lambda(\lambda y_j)$ has the sign of $\sigma Q(y_j)$ for all large $\lambda$.

For large $\lambda$ one also has $x_p<\lambda y_0$. Then $R_\lambda$ alternates at
$x_0<\dots<x_p<\lambda y_0<\dots<\lambda y_q$: the signs alternate within each group, and they differ at $x_p$ and
$\lambda y_0$ by the choice of $\sigma$.
\end{proof}

The same proof applies to sums $\sum_i s_ig_i^{\,r}$ of $r$-th powers of $t$-nomials with signs $s_i=\pm1$, with
$x^{rn}$ in place of $x^{3n}$.

\begin{theorem}\label{thm:four}
Let $\eps=85749603/10^{15}$, $a=98180897/10^{8}$, $b=29608419/10^{9}$ and
\[
g_1=x^{814}-\eps,\quad g_2=-x^{722}+\eps\,x^{1536},\quad g_3=-a\,x^{815}+b\,x^{613},\quad
g_4=a\,x^{721}-b\,x^{923}.
\]
The polynomial $P=g_1^3+g_2^3+g_3^3+g_4^3$ has exactly $15$ positive roots, all simple. Hence $N(4)=15$.
\end{theorem}
\begin{proof}
$P$ has $16$ monomials, so it has at most $15$ positive roots counted with multiplicity. It alternates at the
following sixteen points, with a negative value at the first:
\begin{center}\small
\begin{tabular}{|l|l|l|l|}
\hline
0.9771649 & 0.9809889 & 0.9838065 & 0.9868765\\\hline
0.9902401 & 0.9926549 & 0.9947432 & 0.9978011\\\hline
1.0022038 & 1.0052845 & 1.0073995 & 1.0098561\\\hline
1.0132980 & 1.0164601 & 1.0193795 & 1.0233687
\end{tabular}
\end{center}
The sixteen signs were computed in exact rational arithmetic, and independently in Lean by interval arithmetic with
outward rounding (Appendix~\ref{sec:lean}).
\end{proof}

Here $g_2(x)=-x^{1536}g_1(1/x)$ and $g_4(x)=-x^{1536}g_3(1/x)$, so $x^{4608}P(1/x)=-P(x)$, as in
\eqref{eq:anti}. The example of k-sandhu in the file of \row{251} adds one cube to the three cubes of
Theorem~\ref{thm:three} and has exactly $14$ positive roots: the file gives $14$, and its coefficients have $14$
changes of sign. Appendix~\ref{sec:found} explains how the example above was found.

For five and six cubes the bound \eqref{eq:upper} is already attained by examples in the repository: the file of
\row{253} (junyeopYim) contains twenty points at which a sum of five cubes of binomials alternates, and the file of
\row{255} (k-sandhu) contains twenty-four points for six cubes. The posted statements record only ``more than
$17$'' and ``more than $21$''.

\begin{theorem}\label{thm:main}
$N(1)=1$, $N(2)=3$, and $N(k)=4k-1$ for every $k\ge3$.
\end{theorem}
\begin{proof}
The first two values are Proposition~\ref{prop:two}, and $N(k)\le4k-1$ is \eqref{eq:upper}. Let $a(k)$ be the
largest $n$ such that some polynomial in $\mathcal{S}_k$ alternates at $n+1$ points of $(0,\infty)$; then
$N(k)\ge a(k)$. Theorem~\ref{thm:three}, Theorem~\ref{thm:four} and the file of \row{253} give $a(3)\ge11$,
$a(4)\ge15$ and $a(5)\ge19$. If $k\ge6$ and $a(k-3)\ge4(k-3)-1$, Lemma~\ref{lem:glue} gives
$a(k)\ge a(3)+a(k-3)+1\ge4k-1$. The claim follows by induction on $k$.
\end{proof}

The lower bounds for $k=3,4,5$ are checked in Lean. Lemma~\ref{lem:glue}, the induction and the use of the rule of
signs are not. Without Theorem~\ref{thm:four} the same argument, starting from $k=3,5,6$, gives $N(k)=4k-1$ for all
$k\ge3$ except $k=4$ and $k=7$.

\begin{example}\label{ex:five}
A second example for five cubes has much smaller exponents than the one in the file of \row{253}, whose exponents
have nine digits. Let $\alpha=45239/(5\cdot10^{8})$, $\beta=554277/10^{6}$, $\gamma=170373/500000$,
$\delta=101867/200000$ and
\begin{gather*}
g_1=-\alpha x^{384}+\beta x^{344},\quad g_2=\alpha-\beta x^{40},\quad g_3=\gamma x^{349}-x^{305},\quad
g_4=-\gamma x^{35}+x^{79},\\ g_5=-\delta\,(x^{347}-x^{37}).
\end{gather*}
Then $g_1^3+\dots+g_5^3$ has degree $1152$, twenty monomials and exactly $19$ positive roots. The accompanying
Lean file proves the lower bound $19$.
\end{example}

\begin{corollary}\label{cor:real}
For every $k\ge3$ some polynomial in $\mathcal{S}_k$ has $8k-1$ distinct real roots, and no nonzero polynomial with
at most $4k$ monomials has more.
\end{corollary}
\begin{proof}
If $P=\sum_ig_i^3$ has $4k-1$ positive roots $r_j$, then $x^3P(x^2)=\sum_i\bigl(x\,g_i(x^2)\bigr)^3$ belongs to
$\mathcal{S}_k$ and has the roots $0$ and $\pm\sqrt{r_j}$. Conversely, a nonzero polynomial with at most $4k$
monomials has at most $4k-1$ positive roots, at most $4k-1$ negative roots and possibly the root $0$.
\end{proof}

For $k=3$ this gives $23$ real roots; the file of \row{257} exhibits $22$.

\section{Multiplicity}\label{sec:mult}

Let $M$ be the largest order of a positive root of a nonzero polynomial in $\mathcal{S}_3$. Lemma~\ref{lem:hajos}
gives $M\le11$. What is known is
\[
7\le M\le9 .
\]
The lower bound is the example of maxkieff quoted in the introduction. The upper bound is \row{234}, proved by
Aargau. So for three cubes of binomials the largest order of a root is smaller than the largest number of roots,
which is $11$. For signomials with free coefficients the two numbers are equal.

\subsection{The bound $9$, also for real exponents}

The Lean proof of \row{234} is long, and a short proof does not seem to be known. Its structure is as follows. Let
$r>0$ be a root of order at least $10$ of $P=\sum_is_ig_i^3$. Then the ten numbers
$\sum_k p_k\,r^k\,(k-z)^n$, $n=0,\dots,9$, vanish for every real $z$, where $p_k$ are the coefficients of $P$. If
one of the three inputs is a monomial, $P$ has at most nine monomials and Lemma~\ref{lem:hajos} applies. Otherwise
the proof distinguishes how many of the values $g_i(r)$ vanish; exactly two cannot, because $P(r)=0$. If all three
vanish, the ten equations are shown to force the sums to vanish for all $n$, so that $P=0$. If exactly one
vanishes, the equations are contradictory. If none vanishes, each cube is described by three normalized
parameters, a mean, a variance and a third parameter; either the three means coincide and again the sums vanish
for all $n$, or the equations are shown to be incompatible by two polynomial identities with integer
coefficients, which fill most of the file and are checked by the Lean kernel.

In this argument the exponents enter only as real numbers. This gives the following.

\begin{proposition}\label{prop:real}
Let $G(x)=\sum_{i=1}^{3}c_i\bigl(A_ix^{a_i}+B_ix^{b_i}\bigr)^3$ with real coefficients and real exponents, not
identically zero on $(0,\infty)$. Every positive root of $G$ has order at most $9$.
\end{proposition}
\begin{proof}
After rescaling $x$ the root is $x=1$; put $x=e^{u}$. If $c_iA_iB_i(a_i-b_i)=0$ for some $i$, the $i$-th summand
has at most one term, $G$ has at most nine terms, and the order is at most $8$ by Lemma~\ref{lem:hajos}. Assume
this is not the case. For real $t$ the $n$-th derivative of $e^{-3tu}G(e^{u})$ at $u=0$ is
\[
T_n(t)=\sum_{i=1}^{3}c_i\Bigl[A_i^3\bigl(3a_i-3t\bigr)^n+3A_i^2B_i\bigl(2a_i+b_i-3t\bigr)^n
+3A_iB_i^2\bigl(a_i+2b_i-3t\bigr)^n+B_i^3\bigl(3b_i-3t\bigr)^n\Bigr].
\]
If the root has order at least $10$, then $T_n(t)=0$ for $n\le9$ and all $t$. The case analysis described above,
carried out with real exponents, shows that then $T_n(t_0)=0$ for some $t_0$ and all $n\ge0$. This step is checked
in the Lean file \lean{RealExponentMoments.lean}, which imports the file of \row{234} and uses its lemmas, whose
parameters are real numbers. Now $e^{-3t_0u}G(e^{u})=\sum_e\gamma_e\,e^{(e-3t_0)u}$, where $e$ runs over the
distinct exponents of $G$, and $\sum_e\gamma_e(e-3t_0)^n=0$ for all $n$ forces all $\gamma_e$ to vanish, since
a Vandermonde matrix with distinct nodes is invertible. This contradicts $G\ne0$.
\end{proof}

The algebra in this proof is Aargau's. That it applies to real exponents was observed during the review of this
note, which also produced the Lean file; the file was then recompiled independently. The two steps proved on paper
are the passage from order at least ten to the ten equations $T_n(t)=0$, and the Vandermonde deduction from the
vanishing of all $T_n(t_0)$.

\subsection{Order $9$ with real exponents}

\begin{proposition}\label{prop:nine}
For real $s$,
\[
\Phi(s)=3\sqrt3\,\cosh(s)\,\sinh^3(\sqrt2\,s)\;-\;2\sqrt2\,\sinh^3(\sqrt3\,s)\;=\;\frac{4\sqrt6}{315}\,s^9+O(s^{11}).
\]
Consequently, with $x=e^{s}$, the function
\[
\frac{3\sqrt3}{16}\Bigl[\bigl(x^{\frac13+\sqrt2}-x^{\frac13-\sqrt2}\bigr)^3+\bigl(x^{-\frac13+\sqrt2}-x^{-\frac13-\sqrt2}\bigr)^3\Bigr]
-\frac{\sqrt2}{4}\bigl(x^{\sqrt3}-x^{-\sqrt3}\bigr)^3
\]
is a sum of three cubes of two-term signomials with twelve distinct exponents, and it has a root of order
exactly $9$ at $x=1$. By Proposition~\ref{prop:real}, $9$ is the largest order of a positive root of a nonzero sum
of three cubes of two-term signomials.
\end{proposition}
\begin{proof}
By $\sinh^3u=\tfrac14(\sinh3u-3\sinh u)$ and
$2\cosh s\,\sinh(\lambda s)=\sinh((\lambda+1)s)+\sinh((\lambda-1)s)$, one has
$\Phi=\tfrac{3\sqrt3}{8}F-\tfrac{\sqrt2}{2}G$ with
\begin{align*}
F(s)&=\sinh((3\sqrt2+1)s)+\sinh((3\sqrt2-1)s)-3\sinh((\sqrt2+1)s)-3\sinh((\sqrt2-1)s),\\
G(s)&=\sinh(3\sqrt3\,s)-3\sinh(\sqrt3\,s) .
\end{align*}
For odd $n$, the coefficient of $s^{n}/n!$ in $\tfrac{3\sqrt3}{8}F$ and in $\tfrac{\sqrt2}{2}G$ is $\sqrt6$ times the
following integer.
\begin{center}\small
\begin{tabular}{|l|r|r|r|r|r|}
\hline
\rowcolor{headgray}$n$ & $1$ & $3$ & $5$ & $7$ & $9$\\\hline
$\tfrac{3\sqrt3}{8}F$ & $0$ & $36$ & $1080$ & $29484$ & $801648$\\\hline
$\tfrac{\sqrt2}{2}G$ & $0$ & $36$ & $1080$ & $29484$ & $797040$\\\hline
\end{tabular}
\end{center}
For instance, for $n=3$ the function $F$ gives $(3\sqrt2+1)^3+(3\sqrt2-1)^3-3(\sqrt2+1)^3-3(\sqrt2-1)^3=96\sqrt2$,
and $G$ gives $(3\sqrt3)^3-3(\sqrt3)^3=72\sqrt3$. The difference at $n=9$ is $4608\sqrt6=9!\cdot4\sqrt6/315$. The
second statement follows from
$\cosh s\,\sinh^3(\sqrt2s)=\tfrac12\bigl[(e^{s/3}\sinh(\sqrt2s))^3+(e^{-s/3}\sinh(\sqrt2s))^3\bigr]$.
\end{proof}

The identity was found through a family that goes back to Discussion \#21999. Consider signomials
\begin{equation}\label{eq:common}
G(x)=\sum_{i=0}^{2}w_i\,\bigl(x^{a_i}-x^{b_i}\bigr)^3 ,\qquad a_i<b_i\ \text{real},
\end{equation}
in which all three binomials vanish at $x=1$, and put $L_i=a_i+b_i$, $D_i=b_i-a_i$, $\omega_i=w_iD_i^3$. With
$x=e^t$ the order of $G$ at $x=1$ is the least $n$ such that
\[
\mu_n=\sum_{i=0}^{2}\frac{\omega_i}{D_i^{3}}\sum_{j=0}^{3}\binom3j(-1)^j\bigl(3a_i+jD_i\bigr)^n
\]
is not zero. maxkieff shows that for three distinct values $L_i$ the equations $\mu_3=\dots=\mu_7=0$ determine the
weights $\omega_i$ up to a common factor and the three numbers $D_i^2$, as in the proposition below. This is
part~(i), including all three formulas for the $D_i^2$; the third follows by substituting one of his equations into
another. Parts (ii) and (iii) evaluate the next two moments.

\begin{proposition}\label{prop:family}
Let $p,q$ be real with $pq(p+q)\ne0$, and let $(L_0,L_1,L_2)=(q,-p,0)$, $(\omega_0,\omega_1,\omega_2)=(p,q,-p-q)$,
\[
D_0^2=8p^2+5pq+5q^2,\qquad D_1^2=5p^2+5pq+8q^2,\qquad D_2^2=8p^2+11pq+8q^2 .
\]
These three quadratic forms are positive definite, and
\begin{itemize}\setlength{\itemsep}{0pt}
\item[(i)] $\mu_n=0$ for $n\le7$;
\item[(ii)] $\mu_8=378\,pq\,(p-q)(p+q)(p+2q)(2p+q)$;
\item[(iii)] $\mu_9=-729\,pq\,(p+q)\bigl(6p^4+5p^3q-10p^2q^2+5pq^3+6q^4\bigr)$.
\end{itemize}
Hence $G$ has order exactly $8$ at $x=1$ if $(p-q)(p+2q)(2p+q)\ne0$, and order exactly $9$ otherwise.
\end{proposition}
\begin{proof}
Each $\mu_n$ is a polynomial in $p$, $q$ and the $D_i^2$, since the inner sum is divisible by $D_i^3$ and the
quotient is even in $D_i$. Thus (i) to (iii) are polynomial identities in $p$ and $q$; they were verified with a
computer algebra system. The quartic factor in (iii) equals $12p^4$ for $q=p$ and for $q=-2p$, and $12q^4$ for
$p=-2q$, so $\mu_9\ne0$ on the three exceptional lines.
\end{proof}

For $p=q=1$ one has $L=(1,-1,0)$, $D^2=(18,18,27)$, $\omega=(1,1,-2)$, and
$G(e^{2s/3})=-\tfrac{8}{81\sqrt6}\,\Phi(s)$ with the function $\Phi$ of Proposition~\ref{prop:nine}. The two other
exceptional lines give the same example after a common translation of the $L_i$, a rescaling of the variable, a
permutation of the summands and multiplication by a nonzero constant.

Since orders $8$ and $9$ occur for real exponents, a proof that $M\le8$ or $M=7$ for polynomials cannot be valid
for arbitrary real exponents; it has to use an arithmetic property of the exponents. This is how maxkieff's
argument for \eqref{eq:common} works. By his result quoted above, an example \eqref{eq:common} of order at least
$8$ with rational exponents and distinct $L_i$ would give nonzero rational $p,q$ such that $8p^2+5pq+5q^2$ and
$5p^2+5pq+8q^2$ are squares of rational numbers. After clearing denominators these forms are congruent to $3p^2$
and $3q^2$ modulo $5$, and $3$ is not a square modulo $5$, so $p$, $q$ and the two square roots are divisible by
$5$, which leads to an infinite descent. The reduction of the general case of \eqref{eq:common} to this family is
maxkieff's and was not re-derived here.

\subsection{Integer exponents up to $32$}\label{sec:search}

When not all three binomials vanish at the root, nothing is known that excludes $M=8$ or $M=9$. The following
computation covers small exponents.

\begin{proposition}\label{prop:search}
Let $g_0,g_1,g_2$ be real polynomials with at most two monomials each and all exponents in $[0,32]$. If
$g_0^3+g_1^3+g_2^3\ne0$, every positive root of it has order at most $7$.
\end{proposition}

This was established by computer, as follows; it is not a Lean proof.

\emph{Reductions.} After rescaling $x$ the root is $x=1$. Subtracting the smallest exponent $m$ of the three inputs
from all their exponents divides the sum by $x^{3m}$ and does not change the order of a positive root, so one may
assume that the smallest exponent is $0$. A nonzero sum with a root of order at least $8$ has at least nine
distinct exponents by Lemma~\ref{lem:hajos}. Hence at least two inputs have two monomials, since otherwise the sum
has at most six. An input with one monomial, or the zero polynomial, is written as
$Ax^{a}+Bx^{b}$ with $a<b$ and $A$ or $B$ equal to zero. Two inputs with the same pair of exponents give at most
eight exponents, so the three pairs $(a_i,b_i)$ may be assumed distinct.

\emph{Equations.} Let $g_i=A_ix^{a_i}+B_ix^{b_i}$ with integers $0\le a_i<b_i\le32$. The cube $g_i^3$ has the
exponents $y_{i,j}=3a_i+j(b_i-a_i)$ and coefficients $h_{i,j}=\binom3jA_i^{3-j}B_i^{j}$, $j=0,\dots,3$. A real
quadruple $(h_0,h_1,h_2,h_3)$ has this form for some real $A,B$ if and only if
\begin{equation}\label{eq:cone}
h_1^2=3h_0h_2,\qquad h_2^2=3h_1h_3,\qquad h_1h_2=9h_0h_3 ,
\end{equation}
and the sum has a root of order at least $8$ at $1$ if and only if $\sum_{i,j}h_{i,j}\,y_{i,j}^{\,n}=0$ for
$n=0,\dots,7$. Since nine distinct exponents are needed, the number of coincidences among the twelve numbers
$y_{i,j}$ is at most three.

\emph{The main case.} Fix the exponents, and suppose that $g_P$ has two monomials; normalize $A_P=1$ and write
$z=B_P\ne0$. Suppose that the eight exponents of the two other cubes are distinct. Then the eight linear equations
determine their eight coefficients as cubic polynomials in $z$ with rational coefficients, by Lagrange
interpolation, and \eqref{eq:cone} for the two cubes becomes a system of six polynomial equations of degree at most
$6$ in $z$. A solution exists if and only if these six polynomials have a common real root $z\ne0$. The program
first computes a greatest common divisor of two of them modulo a prime that does not divide a leading coefficient;
if it is constant there is no common root. Otherwise it computes the greatest common divisor of all six over the
rationals and looks for a real root.

Call two cubes \emph{linked} if they share an exponent. If no two cubes are linked, the procedure is applied with
$P=0$ and with $P=1$. If exactly one pair is linked, it is applied with each cube of that pair as $P$. In both
cases two choices of $P$ are tried, and since at least two of the three inputs have two monomials, one of the
choices is an input with two monomials. If exactly two pairs are linked, the procedure is applied with their common
cube as $P$; in the remaining case this cube is a monomial, the sum has exactly nine exponents $e_1,\dots,e_9$,
its coefficient vector is proportional to $\bigl(\prod_{i\ne j}(e_j-e_i)^{-1}\bigr)_j$, and \eqref{eq:cone} is tested
on these numbers.

\emph{Three linked pairs.} Then there are at least two choices of $P$ for which the two other cubes share exactly one
exponent. Indeed, there are at most three coincidences, each pair of cubes accounts for at least one, and an
exponent common to all three cubes accounts for two; so at most one pair of cubes shares two exponents. As before,
one of these choices of $P$ is an input with two monomials. For such $P$ the two other cubes have seven distinct
exponents $E$. The equations $n=0,\dots,6$ determine the seven coefficients of their sum as cubic polynomials in
$z$, and the equation $n=7$ becomes $\sum_j\binom3jW(y_{P,j})z^j=0$ with $W(y)=\prod_{e\in E}(y-e)$. For each real
root $z\ne0$ of this cubic the coefficients of the two cubes are known except for the way the coefficient at the
shared exponent splits between them. The roots are isolated exactly, and the equations \eqref{eq:cone} are then
evaluated in interval arithmetic with rigorous error bounds. The program rejects a candidate when an equation that
does not involve the split, or an equation evaluated at the split determined by another one, has an enclosure
that excludes $0$. It reports a candidate that it can neither reject nor determine. None was reported.

\emph{Result.} For exponents in $[0,32]$ there are $4\,179\,296$ unordered triples of distinct exponent pairs with
smallest exponent $0$; this is also the number of triples in maxkieff's search. Table~\ref{tab:search} gives their
distribution. No triple admits a solution. The programs, their complete output and the versions of the software
are included with this note. During the review of this note the whole computation was repeated with an independent
program that works modulo a prime and uses neither the rational greatest common divisors nor the interval
arithmetic above; it gave the same counts and no solution. The first program was also run in floating-point arithmetic on the
real-exponent examples of Proposition~\ref{prop:family} with $(p,q)=(1,2)$ and $(1,1)$, where the six polynomials
vanish at $z=-1$ to the working precision, as they must. This tests the formulas, not the exact procedure.

\begin{table}[ht]\centering\small
\begin{tabular}{|l|r|}
\hline
fewer than nine distinct exponents (nothing to check) & $3\,787$\\\hline
no linked pair & $2\,145\,738$\\\hline
one linked pair & $1\,673\,517$\\\hline
two linked pairs & $294\,861$\\\hline
three linked pairs & $61\,393$\\\hline
total & $4\,179\,296$\\\hline
\end{tabular}
\caption{The exponent triples of the computation behind Proposition~\ref{prop:search}.}\label{tab:search}
\end{table}

\section{Products of sparse polynomials}\label{sec:products}

If $f$ and $g$ are $t$-nomials and $fg+1$ is not the zero polynomial, it has at most $t^2+1$ monomials and hence
at most $t^2$ positive roots. It is an open question, attributed to Chattopadhyay in \cite{GKPS}, whether the true
maximum is linear in $t$; in 2020 no upper bound of smaller order than $t^2$ and no family of examples with more
than linearly many roots was known \cite{JPSZ}. The files of \row{261}, \row{263} and \row{265} give $8$, $11$ and
$12$ positive roots for $t=3,4,5$, against the bounds $9$, $16$ and $25$. The example of hedanbaomi for trinomials
is
\begin{align*}
f&=x^{46227877}-0.043365112862\,x^{448900690055}+0.001397048929\,x^{987294467927},\\
g&=-1.000331817298-0.043584165131\,x^{454650486655}+0.000109442842\,x^{1454650486655}.
\end{align*}
The ten coefficients of $fg+1$ have the signs $+-+-+--+-+$ in order of increasing degree, with eight changes, so
this polynomial has exactly eight positive roots. All of them lie within $3\cdot10^{-12}$ of $1$; in the variable
$10^{12}\log x$ this is an exponential sum with exponents of size $1$ which is close to the degenerate case $f=1$,
$g=-1$. The method of proof of \row{261}, interval arithmetic in fixed-point numbers with a sign evaluation at nine
points, is reused in most of the other files. Whether $fg+1$ can have nine positive roots for trinomials $f,g$ is
open. A numerical search for such an example, with the method of Appendix~\ref{sec:found}, did not find one and
does not support a conclusion either way.

If $fg+1$ has $r$ distinct positive roots, then $f(x^2)g(x^2)+1$ has at least $2r$ distinct real roots, and
$f(x^2)$, $g(x^2)$ have as many monomials as $f$, $g$. So the three files just mentioned also give at least $16$,
$22$ and $24$ real roots for $t=3,4,5$. The files of \row{267}, \row{269} and \row{271} give $8$, $13$ and $15$ real roots by other
examples. That of \row{267} (BastianMIllan) is $f=x^4-65x^2$, $g=\tfrac{1}{50176}(x^4-65x^2+848)$, for which
$fg+1=\tfrac1{50176}(x^2-1)(x^2-16)(x^2-49)(x^2-64)$. For $pq-rs$ with four $4$-nomials there are at most $32$
monomials, hence at most $31$ positive roots if the polynomial is not zero, and the file of \row{273}
(SemyonOsadchii) exhibits $23$.

\section{Context}\label{sec:context}

Koiran's real $\tau$-conjecture \cite[Conjecture~3]{Koiran} states that a nonzero polynomial
$\sum_{i=1}^k\prod_{j=1}^m f_{ij}$ with real $t$-nomials $f_{ij}$ has at most polynomially many real roots in
$kmt$. Koiran shows that it implies that the permanent has no constant-free arithmetic circuits of polynomial size
and polynomially bounded formal degree. Expanding the products and applying the rule of signs gives the bound
$2kt^m-1$, which is polynomial for fixed $m$. The conjecture is therefore about growing $m$; it is open already for
$k=2$ \cite{Koiran,GKPS,Burgisser}. In the families of this note $m\le3$, and the results have no bearing on
whether the conjecture is true. They determine exact values in the smallest cases. Their content is that for cubes
of binomials the bound from the number of monomials is attained as soon as there are three summands, although it
is far from attained for two (Proposition~\ref{prop:two}), and that it is not attained for multiplicities
(\row{234}).

Bounds for sums of powers of sparse polynomials that do not depend on the exponents of the powers were proved by
Grenet, Koiran, Portier and Strozecki \cite{GKPS}, and by Koiran, Portier and Tavenas \cite{KPT} using Wronskians
as in \cite{VvdP}. For three cubes of binomials the bound of \cite[Theorem~12]{KPT} is about $8\cdot10^{12}$, far
above the eleven given by expansion. Since it does not depend on the power, it becomes smaller than the expansion
bound $3r+2$ for sums of three $r$-th powers of binomials when $r$ is very large. Dutta, Saxena and Thierauf
\cite{DST} relate representations of univariate polynomials as sums of few squares or cubes of sparse polynomials
to circuit lower bounds and identity testing, and Dutta \cite{Dutta} proves the bound for two cubes quoted after
Proposition~\ref{prop:two} and formulates conjectures for more summands. When all the binomials have the same two
exponents, as for sums of cubes of linear forms, the sum has at most four monomials and there is nothing to
prove. Classical results on sums of powers of linear forms, such as Sylvester's theorem \cite{Reznick}, concern
that situation and do not apply to binomials with different exponents.

For multiplicities, Hrube\v{s} \cite{Hrubes} shows that the real $\tau$-conjecture is equivalent to a version that
counts nonzero real roots with multiplicity, and Koiran and Skomra \cite{KoiranSkomra} state a $\tau$-conjecture for
multiplicities and ask for the order of a root of $fg+1$. No published upper bound for the order of a root of a sum
of three cubes of binomials that improves on Lemma~\ref{lem:hajos} was found. Bihan, Dickenstein and Forsg\aa rd
\cite{BDF} ask how the largest number of positive solutions and the largest local multiplicity are related for
sparse systems with free coefficients; the family $\mathcal{S}_3$, whose coefficients are not free, has $11$
against at most $9$. An exact result of a similar kind for another three-term family is due to Bihan and El Hilany
\cite{BihanElHilany}: a plane curve defined by a trinomial meets a line in at most $11$ real points when the
intersection is finite, and $11$ occurs.

The functions $x^{e}$ with distinct real exponents form a Chebyshev system on $(0,\infty)$, which is the reason for
Lemma~\ref{lem:descartes}; see \cite{KarlinStudden,PolyaSzego}. That theory describes the signomials with the
maximal number of roots, but it does not say whether one of them is a sum of cubes of binomials, which is a
nonlinear condition \eqref{eq:cone} on the coefficients.

On 26 November 2026 arXiv, Crossref, zbMATH, MathOverflow, the citation lists of \cite{Koiran,KPT,Dutta} and the
web were searched for real roots of sums of powers of sparse polynomials, the sharpness of the rule of signs in
structured families, the $fg+1$ question, multiplicity bounds, and the specific polynomials above. No statement of
Theorem~\ref{thm:three}, Theorem~\ref{thm:four}, Theorem~\ref{thm:main}, the bound of \row{234} or
Proposition~\ref{prop:nine} was found. The sources
\cite{AlbouyFu,BDF,BihanElHilany,Burgisser,Dutta,DST,GKPS,Hrubes,JPSZ,Koiran,KPT,KoiranSkomra,Lenstra,Reznick}
were read; \cite{Hajos}, the 2026 journal version of \cite{Dutta}, and \cite{KarlinStudden,PolyaSzego,VvdP} were
not, and are cited as standard references. This describes a search on that date and is not a claim of priority.

The project's records describe the eighteen statements as small algebraic cases and do not mention the real
$\tau$-conjecture; the connection made here is the reading of this note, and it is also made in notes of maxkieff.

\section{Open questions}\label{sec:open}

\begin{enumerate}\setlength{\itemsep}{1pt}
\item Can a nonzero sum of three cubes of binomial polynomials have a positive root of order $8$ or $9$?
\item Can $fg+1$ have nine positive roots for trinomials $f,g$?
\item For three cubes of trinomials the known examples have $15$ positive roots and the rule of signs allows $29$.
Lemma~\ref{lem:glue} gives only superadditivity in $k$. Is the bound from the number of monomials attained for
cubes of $t$-nomials with $t\ge3$?
\item For which $r$ does a signed sum of three $r$-th powers of binomials attain $3r+2$ positive roots?
\end{enumerate}

\appendix

\section{The Lean statements and their verification}\label{sec:lean}

Each statement of Table~\ref{tab:rows} is the declaration \lean{PEqualNP.Outside.L\_gct\_N}. Write
\begin{center}
$\Sigma_k$ for \lean{(Finset.univ : Finset (Fin k)).sum fun i => Polynomial.C (s i) * g i \^{} 3},\\[2pt]
$\mathrm{pos}(P)$ for \lean{(P.roots.toFinset.filter (fun x : ℝ => 0 < x)).card},\qquad
$\mathrm{all}(P)$ for \lean{P.roots.toFinset.card},
\end{center}
and $H_t$ for the hypotheses \lean{∀ i, s i = 1 ∨ s i = -1} and \lean{∀ i, (g i).support.card ≤ t}. In Mathlib the
roots of the zero polynomial form the empty multiset, so the zero polynomial has no roots in these statements; the
hypothesis \lean{support.card ≤ t} allows fewer than $t$ monomials.

\begin{center}\footnotesize
\begin{longtable}{|L{2.2cm}|L{7.4cm}|L{4.7cm}|}
\hline
\rowcolor{headgray}\textbf{name} & \textbf{statement} & \textbf{declaration that counts roots}\\\hline
\endhead
\row{234} & $\forall s\,g,\ H_2\to\Sigma_3\ne0\to\forall x,\ 0<x\to{}$\lean{rootMultiplicity x} $\Sigma_3\le9$ &
\lean{SparseCubes.multiplicity\_bound}\\\hline
\row{241} & $\neg\,(\forall s\,g,\ H_2\to\mathrm{pos}(\Sigma_3)\le9)$ & \lean{eleven\_pos\_roots}\\\hline
\row{251}, 253, 255 & $\neg\,(\forall s\,g,\ H_2\to\mathrm{pos}(\Sigma_k)\le B)$ with
$(k,B)=(4,13),(5,17),(6,21)$ & \lean{roots\_lower}, with \lean{count} $=14,19,23$\\\hline
\row{243}, 245, 247, 249 & $\neg\,(\forall s\,g,\ H_t\to\mathrm{pos}(\Sigma_3)\le B)$ with
$(t,B)=(3,14),(4,21),(5,26),(6,33)$ & \lean{roots\_lower}, with \lean{count} $=15,23,27,35$\\\hline
\row{257}, 259 & $\neg\,(\forall s\,g,\ H_t\to\mathrm{all}(\Sigma_3)\le B)$ with $(t,B)=(2,19),(3,29)$ &
\lean{roots\_even\_lower}, giving $2\cdot11$ and $2\cdot15$\\\hline
\row{261}, 263, 265 & $\neg\,(\forall f\,g,\ $supports${}\le t\to\mathrm{pos}(fg+1)\le B)$ with
$(t,B)=(3,7),(4,10),(5,11)$ & \lean{eightRoots}; \lean{roots\_lower}, with \lean{count} $=11,12$\\\hline
\row{267}, 269, 271 & $\neg\,(\forall f\,g,\ $supports${}\le t\to\mathrm{all}(fg+1)\le B)$ with
$(t,B)=(3,7),(4,12),(5,14)$ & \lean{sub\_roots}; \lean{main}; \lean{exists\_roots}\\\hline
\row{273} & $\neg\,(\forall p\,q\,r\,s,\ $supports${}\le4\to\mathrm{pos}(pq-rs)\le22)$ & \lean{roots\_lower}, with
\lean{count} $=23$\\\hline
\end{longtable}
\end{center}

Most files evaluate the polynomial at rational points by fixed-point interval arithmetic with outward rounding,
executed by the Lean kernel, and apply the intermediate value theorem; this method was introduced in the file of
\row{261} by hedanbaomi. The file of \row{241} evaluates at the algebraic points of Theorem~\ref{thm:three} through
rational enclosures. The file of \row{267} verifies eight explicit integer roots; that of \row{269} combines twelve
sign changes between rational sample points with the root $0$; that of \row{271} uses fifteen sign changes between
consecutive half-integers. The files of \row{257} and \row{259} obtain real roots from positive roots of an even
polynomial.

Apart from \row{234}, the files prove lower bounds. Mathlib contains the rule of signs for polynomials
(\lean{Polynomial.roots\_countP\_pos\_le\_signVariations}), but none of the files uses it, and the upper bounds and
exact values for the number of roots stated in this note are not formalized. A Lean proof of
Theorem~\ref{thm:main} would in addition need Lemma~\ref{lem:glue}.

For this note the eighteen files were compiled with Lean~4.36.0 and the Mathlib version fixed by the repository,
after removing the line \lean{import PEqualNP.Defs} where present, so that only Mathlib is imported. All compile,
and \lean{\#print axioms} lists \lean{propext}, \lean{Classical.choice} and \lean{Quot.sound} for each of the
eighteen theorems. The state of the repository was commit \texttt{ff9577c}.

Three Lean files accompany this note. \lean{FourCubes15.lean} proves
\begin{center}
$\neg\,(\forall s\,g,\ H_2\to\mathrm{pos}(\Sigma_4)\le14)$
\end{center}
with the polynomial of Theorem~\ref{thm:four}, and \lean{FiveCubes19.lean} proves the same with $\Sigma_5$ and $18$
for Example~\ref{ex:five}. Their proofs reuse without change the interval arithmetic and the counting argument of
the file of \row{251} (k-sandhu, after hedanbaomi); only the polynomial, the sample points and the expanded
coefficients are new. Both import Mathlib only. \lean{RealExponentMoments.lean} contains the step of
Proposition~\ref{prop:real}: for real exponents and nonzero coefficients, if the sums $T_n(t)$ vanish for $n\le9$
and all $t$, then for some $t_0$ they vanish for all $n$. It imports the file of \row{234}. The three files compile
with the same versions and depend on the same three axioms. They were not submitted to the project's own checking
service, which accepts posted statements only.

The supplementary directory contains the three Lean files, a program that evaluates the polynomials of
Theorem~\ref{thm:four} and Example~\ref{ex:five} in exact rational arithmetic at the rational points stored in
their certificates (the sixteen points of Theorem~\ref{thm:four} are those printed above), the programs for
Propositions \ref{prop:nine}, \ref{prop:family} and \ref{prop:search} with their output, a script that repeats the
compilation of the Lean files, and a record of the reviews of the note.

\section{How the new examples were found}\label{sec:found}

The examples of Theorem~\ref{thm:four} and Example~\ref{ex:five} were found numerically; their proofs do not depend
on how. A signomial with $n+1$ given exponents and $n$ given simple positive roots is unique up to a constant
factor. For a sum of $k$ cubes of binomials the $4k$ exponents are determined by the $2k$ exponents of the
binomials. Taking these and $4k-1$ roots as unknowns, one computes the coefficient vector by solving a linear
system and asks that the four coefficients of each cube satisfy \eqref{eq:cone}. This is a nonlinear system with
more unknowns than equations, which was solved by least squares from random starting points. The search was
restricted to polynomials with the symmetry \eqref{eq:anti}: the binomials come in pairs $g(x)$,
$-x^{N}g(1/x)$, possibly with binomials $c\,(x^{N/2+m}-x^{N/2-m})$, and then only half of the roots are unknown.
For $k=4$ solutions were found with two pairs, and for $k=5$ with two pairs and one further binomial. A numerical
solution was made exact by rounding the exponents to rational numbers with a common denominator, solving again
with these exponents, rounding the coefficients, and checking the signs in exact arithmetic. Since the roots are
simple, small changes of the parameters do not destroy them.

\section*{Credits}
The statements of Table~\ref{tab:rows} were posted by the PEqualNP project. The accepted proofs were submitted by
Aargau, BastianMIllan, blarneys, gcq, hedanbaomi, junyeopYim, k-sandhu, SemyonOsadchii and XXXLUMI, as listed in the
table; some of these submissions record the assistance of AI systems. An earlier pull request for \row{234} by
williamzheng0711 (\#23279) was closed without being accepted. The example of XXXLUMI in Theorem~\ref{thm:three} is
the starting point of the later files on cubes, and the interval arithmetic of hedanbaomi is used in most of them
and in two files of this note. The example of order $7$, the bound for binomials with a common root, part~(i) of
Proposition~\ref{prop:family} and the descent modulo $5$ are due to maxkieff, and the obstruction for three
consecutive unshared exponents is due to silentoplayz. The algebra behind Proposition~\ref{prop:real} is Aargau's.
The remaining statements of Sections \ref{sec:more} and \ref{sec:mult} are those listed as (a) to (d) in the
introduction.

\begin{thebibliography}{99}\small
\bibitem{AlbouyFu} A.~Albouy and Y.~Fu, Some remarks about Descartes' rule of signs, \emph{Elem.\ Math.}\ 69 (2014)
186--194. arXiv:1307.6842.
\bibitem{BDF} F.~Bihan, A.~Dickenstein and J.~Forsg\aa rd, Sparse systems with high local multiplicity,
\emph{J.\ London Math.\ Soc.}\ (2) 111 (2025), no.~3, e70106. arXiv:2402.08410.
\bibitem{BihanElHilany} F.~Bihan and B.~El Hilany, A sharp bound on the number of real intersection points of a
sparse plane curve with a line, \emph{J.\ Symbolic Comput.}\ 81 (2017) 88--96. arXiv:1506.03309.
\bibitem{Burgisser} P.~B\"urgisser, Completeness classes in algebraic complexity theory, arXiv:2406.06217 (2024).
\bibitem{Dutta} P.~Dutta, Real $\tau$-conjecture for sum-of-squares: a unified approach to lower bound and
derandomization, in: \emph{CSR 2021}, Lecture Notes in Comput.\ Sci., Springer, 2021, pp.~78--101. Journal version:
\emph{Theory Comput.\ Syst.}\ 70 (2026), article 52.
\bibitem{DST} P.~Dutta, N.~Saxena and T.~Thierauf, A largish sum-of-squares implies circuit hardness and
derandomization, in: \emph{ITCS 2021}, LIPIcs 185, 23:1--23:21. Journal version: \emph{comput.\ complexity} 33
(2024), article~3.
\bibitem{GKPS} B.~Grenet, P.~Koiran, N.~Portier and Y.~Strozecki, The limited power of powering: polynomial
identity testing and a depth-four lower bound for the permanent, in: \emph{FSTTCS 2011}, LIPIcs,
doi:10.4230/LIPIcs.FSTTCS.2011.127. arXiv:1107.1434.
\bibitem{Hajos} G.~Haj\'os, [Solution to problem 41] (in Hungarian), \emph{Mat.\ Lapok} 4 (1953) 40--41. Cited after
\cite{Lenstra,KoiranSkomra}.
\bibitem{Hrubes} P.~Hrube\v{s}, On the real $\tau$-conjecture and the distribution of complex roots, \emph{Theory
Comput.}\ 9 (2013) 403--411.
\bibitem{JPSZ} G.~Jindal, A.~Pandey, H.~Shukla and C.~Zisopoulos, How many zeros of a random sparse polynomial are
real?, in: \emph{ISSAC 2020}, 273--280. arXiv:1911.02540.
\bibitem{KarlinStudden} S.~Karlin and W.~J.~Studden, \emph{Tchebycheff Systems: With Applications in Analysis and
Statistics}, Interscience, 1966.
\bibitem{Koiran} P.~Koiran, Shallow circuits with high-powered inputs, in: \emph{Innovations in Computer Science
(ICS 2011)}, Tsinghua University Press, 309--320. arXiv:1004.4960.
\bibitem{KPT} P.~Koiran, N.~Portier and S.~Tavenas, A Wronskian approach to the real $\tau$-conjecture,
\emph{J.\ Symbolic Comput.}\ 68 (2015) 195--214. arXiv:1205.1015.
\bibitem{KoiranSkomra} P.~Koiran and M.~Skomra, Intersection multiplicity of a sparse curve and a low-degree curve,
\emph{J.\ Pure Appl.\ Algebra} 224 (2020) 106279. arXiv:1904.00702.
\bibitem{Lenstra} H.~W.~Lenstra, Jr., Finding small degree factors of lacunary polynomials, in: \emph{Number Theory
in Progress}, Vol.~1 (Zakopane--Ko\'scielisko, 1997), de Gruyter, 1999, 267--276.
\bibitem{PolyaSzego} G.~P\'olya and G.~Szeg\H{o}, \emph{Problems and Theorems in Analysis II}, Springer, 1976, Part
Five.
\bibitem{Reznick} B.~Reznick, On the length of binary forms, in: \emph{Quadratic and Higher Degree Forms},
Dev.\ Math.\ 31, Springer, 2013, 207--232. arXiv:1007.5485.
\bibitem{VvdP} M.~Voorhoeve and A.~J.~van der Poorten, Wronskian determinants and the zeros of certain functions,
\emph{Indag.\ Math.}\ 78 (1975) 417--424.
\end{thebibliography}

\end{document}
